\documentclass[conference,onecolumn]{IEEEtran}

\usepackage[T1]{fontenc}
\usepackage[utf8]{inputenc}
\usepackage[spanish,es-tabla]{babel}
\usepackage{lmodern}
\usepackage{microtype}
\usepackage{csquotes}
\usepackage{amsmath,amssymb}
\usepackage{booktabs}
\usepackage{tabularx}
\usepackage{array}
\usepackage{url}
\usepackage{float}
\usepackage[hidelinks]{hyperref}
\usepackage[backend=biber,style=ieee,sorting=none,isbn=false,url=true,doi=true]{biblatex}
\addbibresource{referencias.bib}

\usepackage{ragged2e}
\newcolumntype{Y}{>{\RaggedRight\arraybackslash}X}
\newcolumntype{L}[1]{>{\RaggedRight\arraybackslash}p{#1}}
\renewcommand{\arraystretch}{1.13}
\newcommand{\si}{\ensuremath{\checkmark}}
\newcommand{\no}{\ensuremath{\times}}
\newcommand{\req}[2]{\par\noindent\textbf{#1:} #2\par\smallskip}
\setlength{\emergencystretch}{1em}
\AtBeginDocument{\spanishplainpercent}
\renewcommand{\IEEEkeywordsname}{Palabras clave}

\begin{document}
\title{Especificación preliminar de requisitos\\\large Sistema de cobro de impuestos prediales en línea para un GAD cantonal del Ecuador}
\author{\IEEEauthorblockN{TA01 - Mendoza Palma Allan Jeremy}
\IEEEauthorblockA{Universidad Técnica Estatal de Quevedo (UTEQ)\\Modalidad individual}}
\maketitle
\vspace{-1.5em}
\noindent\textbf{Base metodológica:} ISO/IEC/IEEE 29148:2018 e ISO/IEC 25010:2023 \cite{iso29148,iso25010}.

\section*{1. Identificación de stakeholders}
\begin{table}[H]
\centering
\caption{Stakeholders y necesidades preliminares}
\label{tab:stakeholders}
\scriptsize
\setlength{\tabcolsep}{3pt}
\begin{tabularx}{\columnwidth}{@{}L{0.28\columnwidth}L{0.26\columnwidth}Y@{}}
\toprule
\textbf{Stakeholder} & \textbf{Rol} & \textbf{Necesidades}\\
\midrule
Contribuyentes & Usuarios que consultan y cancelan obligaciones. & Consultar saldo por clave catastral; obtener constancia de pago.\\
Dirección Financiera / Tesorería del GAD & Recaudación y control de ingresos. & Aplicar reglas tributarias aprobadas; rastrear pagos confirmados.\\
Dirección de Avalúos y Catastros & Administración de información predial. & Identificar claves catastrales; proporcionar bases imponibles actualizadas.\\
Servicio de Rentas Internas (SRI) & Intercambio institucional autorizado. & Recibir datos acordados; validar la estructura de los intercambios.\\
Proveedor de pasarela local & Procesamiento de transacciones. & Recibir solicitudes conforme a contrato de interfaz; comunicar estados verificables.\\
\bottomrule
\end{tabularx}
\end{table}


\section*{2. Requisitos funcionales y no funcionales}
\subsection*{2.1. Requisitos funcionales}
\req{RF-01}{El sistema deberá mostrar el saldo pendiente asociado a una clave catastral cuando el contribuyente solicite la consulta con una clave catastral válida, con el fin de identificar el importe adeudado.}
\req{RF-02}{El sistema deberá calcular el total exigible en dólares estadounidenses sumando impuesto predial, intereses y recargos según la tabla tributaria versionada aprobada por Tesorería cuando exista una obligación impaga, con el fin de determinar el importe cobrable.}
\req{RF-03}{El sistema deberá procesar el pago electrónico de una obligación predial mediante la pasarela contratada cuando el contribuyente confirme el importe calculado, con el fin de permitir la recaudación en línea.}
\req{RF-04}{El sistema deberá emitir un comprobante municipal de pago con identificador único, clave catastral, importe, fecha y referencia de transacción cuando la pasarela confirme el pago como aprobado, con el fin de respaldar la operación realizada.}
\req{RF-05}{El sistema deberá conciliar cada registro de recaudación predial intercambiado con el SRI contra su acuse correspondiente cuando el SRI devuelva la respuesta por el canal interinstitucional autorizado, con el fin de identificar registros rechazados.}


\subsection*{2.2. Requisitos no funcionales}
\req{RNF-01}{El sistema deberá responder al menos el 95\,\% de las consultas de deuda predial en un máximo de 3,0 segundos con 300 usuarios virtuales concurrentes cuando se ejecute una carga sostenida de 30 minutos, con el fin de limitar los tiempos de espera en periodos de vencimiento.}
\req{RNF-02}{El sistema deberá presentar 0 hallazgos confirmados de vulnerabilidades altas o críticas con puntuación CVSS v3.1 mayor o igual a 7,0 en las interfaces web y API expuestas cuando se ejecute la prueba de penetración previa a producción, con el fin de disminuir el riesgo de acceso no autorizado.}
\req{RNF-03}{El sistema deberá mantener una disponibilidad mensual de al menos 99,5\,\% en el servicio público de consulta predial cuando se supervise durante todo el mes calendario, con el fin de asegurar continuidad de la atención digital.}
\req{RNF-04}{El sistema deberá permitir que al menos 18 de 20 contribuyentes completen una consulta de deuda en un máximo de 120 segundos sin asistencia cuando utilicen un teléfono móvil con una clave catastral de prueba válida, con el fin de reducir las barreras de uso.}
\req{RNF-05}{El sistema deberá clasificar correctamente 100 de 100 respuestas válidas de la pasarela según su contrato de interfaz cuando se ejecuten pruebas de integración en el ambiente de certificación, con el fin de asegurar la interoperabilidad del cobro electrónico.}


\subsection*{2.3. Verificación de calidad de redacción}
Se comprueban los nueve atributos indicados para el ejercicio, tomando como referencia ISO/IEC/IEEE 29148:2018 \cite{iso29148}.
\begin{table}[H]
\centering
\caption{Evaluación documental de los diez requisitos}
\label{tab:evaluacion}
\scriptsize
\setlength{\tabcolsep}{2.4pt}
\begin{tabular}{@{}l*{9}{c}@{}}
\toprule
\textbf{ID} & \textbf{N} & \textbf{Ap} & \textbf{NA} & \textbf{C} & \textbf{S} & \textbf{F} & \textbf{V} & \textbf{Co} & \textbf{Cf}\\
\midrule
RF-01 & \si & \si & \si & \si & \si & \si & \si & \si & \si\\
RF-02 & \si & \si & \si & \si & \si & \si & \si & \si & \si\\
RF-03 & \si & \si & \si & \si & \si & \si & \si & \si & \si\\
RF-04 & \si & \si & \si & \si & \si & \si & \si & \si & \si\\
RF-05 & \si & \si & \si & \si & \si & \si & \si & \si & \si\\
RNF-01 & \si & \si & \si & \si & \si & \si & \si & \si & \si\\
RNF-02 & \si & \si & \si & \si & \si & \si & \si & \si & \si\\
RNF-03 & \si & \si & \si & \si & \si & \si & \si & \si & \si\\
RNF-04 & \si & \si & \si & \si & \si & \si & \si & \si & \si\\
RNF-05 & \si & \si & \si & \si & \si & \si & \si & \si & \si\\
\bottomrule
\end{tabular}
\end{table}
\noindent\footnotesize\textbf{Leyenda:} N: necesario; Ap: apropiado; NA: no ambiguo; C: completo; S: singular; F: factible; V: verificable; Co: correcto; Cf: conforme; \si: cumplimiento preliminar. La factibilidad y corrección se ratifican con el GAD y contratos de integración, especialmente RF-05.\normalsize


\section*{3. Clasificación según ISO/IEC 25010:2023}
\begin{table}[H]
\centering
\caption{Características de calidad de los RNF \cite{iso25010}}
\label{tab:clasificacion}
\scriptsize
\setlength{\tabcolsep}{3pt}
\begin{tabularx}{\columnwidth}{@{}L{0.13\columnwidth}L{0.34\columnwidth}Y@{}}
\toprule
\textbf{RNF} & \textbf{Característica} & \textbf{Justificación}\\
\midrule
RNF-01 & Eficiencia de desempeño & Mide tiempo de respuesta bajo concurrencia definida.\\
RNF-02 & Seguridad & Limita a cero las vulnerabilidades de severidad alta o crítica.\\
RNF-03 & Fiabilidad & Establece disponibilidad mensual mínima.\\
RNF-04 & Capacidad de interacción & Mide la culminación autónoma de la consulta en un tiempo fijado.\\
RNF-05 & Compatibilidad & Verifica interpretación de respuestas según contrato de interfaz.\\
\bottomrule
\end{tabularx}
\end{table}

\printbibliography[title={Referencias}]
\end{document}
